\section{Calibration of a terminal marked field on its selected history}
\label{sec:actual-marked-portfolio-calibration}

The local invitation in Section~\ref{sec:ex-marked-field-local} is a
terminal-scale construction. The public field calibration also needs one
label correction that works at all earlier scales of the already selected
history. This section follows that exact route: the campaign supplies the
fixed portfolio, the terminal Cream supplies the final coefficient, and a
backward induction through the selected Music relocates that coefficient.
The history, doors, family, and Music remain fixed throughout.

\subsection{From the campaign to the fixed portfolio}

The campaign produced in Section~\ref{sec:marked-campaign-producer} is an
imprint with the powered slot-degree bound. Applying the fixed-stencil
construction of Theorem~\ref{thm:ex-marked-portfolio-section} gives, at
each label degree, one finite section, stencil, and portfolio. These choices
are made before selecting a finite-volume history. On the public branch the
same data are attached to the selected family: at degree $d$ the section is
the family section, and its portfolio certificate is retained on the same
family. In particular, the hypotheses of the local reserve proposition
remain attached to that family and its terminal Keeper.

The precise same-history obligation is this attachment of the fixed
portfolio to the selected family. The subsequent argument does not replace
that history by one chosen for each period, degree, or terminal volume.

\subsection{Finite backward calibration}

Fix one selected history $H$, period $p$, terminal volume $N$, label frame
$F$, action $A$, section $B$, and a Music $M$ on this history. At scale $i$,
write $u_i:F\to F$ for the Music index, $\Pi_i(v)$ for its played coefficient,
and $g_i(v)$ for its charge. A state at scale $i$ consists of a primary
$a_i\in F$, a residual coefficient $Z_i$, and a bill $b_i$. Its translated
state at displacement $v$ is
\[
 \mathcal Z_i(v)=Z_i+\Pi_i(a_i+v),\qquad
 c_i(v)=b_i+g_i(a_i+v).
\]
Call the state keyed if $Z_i.x=0$ and $Z_i$ has the prescribed format and
remainder bound. Write $J_i$ for the history's clock index at scale $i$;
the schedule has $J_{i+1}\leq J_i+1$. The source transition from scale
$i+1$ to $i$ preserves keyedness, gives the exact density push identity
for every $v$, and gives the direction relation comparing the leading
coordinates for every $v$. It does not assert equality of those coordinates.
Writing $L_i$ for the physical side at step $i$,
$e_i=\operatorname{barunit}(L_i)$, and $K_i=W(J_{i+1})$, for each
displacement $v$ there is an invertible map $B_{i,v}$ such that
\[
\begin{gathered}
 \|B_{i,v} x\|_1,\ \|B_{i,v}^{-1}x\|_1\leq K_i\|x\|_1,
 \qquad c_i(v),t_{i+1}(v)\geq0,\\
 \operatorname{core}(B_{i,v}x)=B_{i,v}(\operatorname{core}x),\\
 \|x_i(v)-B_{i,v}(\operatorname{vowel}_{L_i}(x_{i+1}(v)))\|_1
   \leq e_i(\|\operatorname{core}(x_{i+1}(v))\|_1+t_{i+1}(v)),\\
 c_i(v)\leq K_i\|\operatorname{core}(\operatorname{vowel}_{L_i}(x_{i+1}(v)))\|_1
       +e_i t_{i+1}(v).
\end{gathered}
\]
Here $x_i(v),x_{i+1}(v)$ are the leading coordinates of the translated
states, and $t_{i+1}(v)$ is the upper state's nonnegative bill. More
explicitly, these coordinates are the corresponding Music-index values
$u_i(a_i+v)$ and $u_{i+1}(a_{i+1}+v)$; the displayed direction inequality,
not an equality, compares them. The maps may depend on $v$, but the
displayed constants do not. The density identity is
$\rho_i(v)=L_i^d\,\mathsf P_i\rho_{i+1}(v)$, with $\mathsf P_i$ exactly
the transition push in the history. The finite Tale witness supplies the
weighted support estimate for the specified exponential entries and forced
profiles. Every bound is required at every displacement, not only at the
base displacement $v=0$.

For clarity, the paid one-step input has the following content. At scale
$i$ put $C_i=\operatorname{columns}(i)$. The tuned coupon has a number
$\eta_i\geq0$ with $(1+2C_i)\eta_i\leq e_i$. For every keyed upper
coefficient $U$ with zero leading coordinate, format, and bill $q$, it
provides a finite wing and a coefficient correction whose leading error
cost is at most $\eta_i q$, whose remainder cost is at most $\eta_i q$,
and whose support charge is at most $2q$. The associated Carol witness is
a finite wing that works for every displacement and compares the upper and
lower forced profiles, exponential entries, and charges. The one-step
construction takes the correction coefficient $V$, sets
$b=u_i^{-1}V.x$, and uses
\[
 a_i=a_{i+1}+b,\qquad
 Z_i=V+\Pi_i(-b),\qquad
 b_i=\eta_i q+g_i(-b).
\]
The leading coordinate of $Z_i$ vanishes because
$V.x+u_i(-b)=0$. The format and remainder bounds follow by adding the
coupon coefficient to the earned played coefficient. For every $v$, the
linear density readout cancels the two terms at $b$ and $-b$; the remaining
density is exactly the transition push of the upper state's density. The
direction estimate follows from the coupon's paid direction bounds, the
Music cone estimate $g_i(\pm b)\leq C_i\eta_iq$, the Carol charge
comparison, and $(1+2C_i)\eta_i\leq e_i$. Finally, union the two finite
wings and add their Tale inequalities; this gives one common finite support
for all $v$. This is the precise one-step content used in the finite descent.

\begin{lemma}[Backward propagation of one terminal coefficient]
\label{lem:actual-marked-antilles}
Let $0\leq r\leq l\leq m\leq N$. Suppose the history is selected and the
same Music is heard and responds at every level from $r$ through $N$.
Suppose its one-step response on each door $D_i$, $l\leq i<m$, is the
response described above: for every keyed input it supplies a keyed output,
the all-displacement density identity and direction bound, and a finite Tale
witness, with the tuning and Coupon bounds stated there. At level $m$, let $Z$ be a
coefficient with the required format, remainder bound $t$, and
\[
             \|Z.x\|_1+t\leq w.
\]
Assume also the scale room and bank budgets. Then there are keyed states
$(a_i,Z_i,b_i)$ for $l\leq i\leq m$ such that:
\begin{enumerate}
\item the terminal state is the state obtained by adjoining $Z$ to the
  level-$m$ Music and subtracting the Music's primary which represents
  $Z.x$;
\item every adjacent pair is related by the supplied response on $D_i$;
\item put $D_i(v)=Z_i+\Pi_i(a_i+v)$. With $a_l$ the bottom primary,
  set $x_i=D_i(-a_l).x$ and $c_i=b_i+g_i(a_i-a_l)$. Then, for every
  $l\leq i\leq m$,
  \[
    \exp\!\left(\frac{\sigma_m-\sigma_i}{8}\right)
       (c_i+\|x_i\|_1)\leq 4(1+C_m)w,
    \qquad C_m=\operatorname{columns}(m),
  \]
  where $\sigma_i$ is the history scale parameter.
\end{enumerate}
The terminal state has, for every $v\in F$, exact identities
\[
 \mathcal Z_m(v).x=Z.x+u_m(v),\qquad
 \mathcal Z_m(v).\mathrm{density}
       =Z.\mathrm{density}+\Pi_m(v).\mathrm{density},
\]
and its terminal bill obeys
\[
 c_m(v)\leq w+C_m\bigl(w+\|\mathcal Z_m(v).x\|_1\bigr),
       \qquad C_m=\operatorname{columns}(m).
\]
\end{lemma}
\begin{proof}
At the terminal level set $a=u_m^{-1}(Z.x)$. In the public application
$m=N$, and the top clause of `Sung` gives
$(M.\mathrm{tone}\,N).\mathrm{index}=\mathrm{id}$: its played coefficient
has leading value $v$ at every label $v$. Thus in that application
$u_N(v)=v$. Form the residual
coefficient by adding to $Z$ the played field $\Pi_m(-a)$, and set its bill
to $t+g_m(-a)$. Its leading coordinate is zero because
$\Pi_m(-a).x=u_m(-a)=-Z.x$. This proves keyedness. For displacement $v$,
linearity of the Music readout gives
\[
 (Z+\Pi_m(-a)+\Pi_m(a+v)).x=Z.x+u_m(v),
\]
and linearity of its density gives cancellation of the $-a$ and $a$
terms, leaving $Z$ plus the density of $\Pi_m(v)$. These are exact
identities.

For the charge bound, the Music's cone estimate gives
$g_m(x)\leq C_m\|u_m(x)\|_1$. Since $u_m(a)=Z.x$ and
$u_m(a+v)=\mathcal Z_m(v).x$, it follows that
\[
 c_m(v)=t+g_m(-a)+g_m(a+v)
 \leq t+C_m\|Z.x\|_1+C_m\|\mathcal Z_m(v).x\|_1
 \leq w+C_m\bigl(w+\|\mathcal Z_m(v).x\|_1\bigr).
\]
This is the displayed terminal bill bound.

We now descend from $m$ to $l$. At a step $i$, assume a keyed state at
$i+1$ has been constructed. The same Music supplies its Heard and
Responds data at $i$ and $i+1$. Apply the paid one-step response with the
state's format and remainder bound; its zero leading coordinate is exactly
the Coupon hypothesis. The output is a keyed state at $i$, its Music-index
formula, the density push identity, the paid direction inequality, and the
finite forced-profile witness. Replace only the state
at $i$ and keep all states at levels above $i$ unchanged. Thus all already
proved chair and transition clauses persist, and the new clauses hold at
$i$. Finite downward induction constructs the full sequence. This is a
finite induction on $m-l$; it does not pass to a subsequence or alter the
selected history.

For the signed difference, let $a_l$ be the primary of the bottom state
and evaluate every state at displacement $-a_l$. The bottom played
coordinate is then zero. At each level, the format and remainder clauses
follow from keyedness and the Music's earned bounds; transition and density
clauses follow from the all-displacement response identities. For the
numerical propagation fix this displacement $v=-a_l$ and abbreviate
$x_i=x_i(v)$, $c_i=c_i(v)$, and $t_i=t_i(v)$; thus the map used at step $i$
is the possibly displacement-dependent $B_{i,v}$, whose norm bounds are
uniform. Put $C_i=\operatorname{columns}(i)$,
$s_i=L_i^{1/8}$, $e_i=s_i^{-3}$, and $A_i=32(1+C_i)^2$. Room cost says
$64W(J_i+1)^2\leq C_i$; since $J_{i+1}\leq J_i+1$ and $W$ is
nondecreasing, this yields $64K_i^2\leq C_i$. The bank inequality at
step $i$, together with $W(J_{i+1})\leq W(J_{i+1}+1)$, yields
$100000K_i^6(1+C_{i+1})^2\leq s_i$. Also $K_i\geq1$ and
$C_i\geq1$. Consequently
\[
 4K_ie_i\leq1,\qquad 4K_iA_{i+1}e_i\leq1,
       \qquad (4K_i^2+1)(s_ie_i)\leq1.
\]
Indeed, writing $C=C_{i+1}$ and $K=K_i$, the bank bound implies
$s_i\geq128K^2(1+C)^2$, since $100000K^6\geq128K^2$ for $K\geq1$.
Thus $4KA_{i+1}=128K(1+C)^2\leq s_i$ and $4K\leq s_i$.
Also $(4K^2+1)s_ie_i=(4K^2+1)s_i^{-2}\leq1$ because
$s_i\geq128K^2$. These are the three displayed numerical inequalities.
The direction inequalities for the translated states imply, with
$u=\|x_i\|_1$, $v=\|x_{i+1}\|_1$, and
$y=\|\operatorname{vowel}_{L_i}(x_{i+1})\|_1$,
\[
 y\leq K_i(u+e_i(y+c_{i+1})),\qquad
 c_i\leq K_i y+e_i c_{i+1},\qquad v\leq y.
\]
If $A_i u\leq c_i$, combining these inequalities with the three
numerical bounds gives
$A_{i+1}v\leq c_{i+1}$ and $s_i c_i\leq c_{i+1}$.
Indeed, substituting $u\leq(K_i y+e_i c_{i+1})/A_i$ into the first
direction inequality gives
\[
 y\leq \left(\frac{K_i^2}{A_i}+K_ie_i\right)y
       +\left(\frac{K_ie_i}{A_i}+K_ie_i\right)c_{i+1}.
\]
Since $A_i=32(1+C_i)^2$, we have $K_i^2/A_i\leq1/32$;
$K_ie_i\leq1/4$; and $A_i\geq1$. Thus the coefficient of $y$ is at
most $9/32<1/2$, and the coefficient of $c_{i+1}$ is at most
$2K_ie_i$. It follows that $y\leq4K_ie_i c_{i+1}$. Now $v\leq y$ and
$4K_iA_{i+1}e_i\leq1$ give $A_{i+1}v\leq c_{i+1}$. Also
$c_i\leq K_i y+e_i c_{i+1}$ gives
$s_ic_i\leq(4K_i^2+1)s_ie_i c_{i+1}\leq c_{i+1}$ by the third
numerical bound. At the bottom,
$keyedness$ gives $x_l=0$ and nonnegative cost. Thus induction gives
$A_i\|x_i\|_1\leq c_i$ at every level, and iterating
$s_i c_i\leq c_{i+1}$ gives
$\prod_{j=i}^{m-1}s_j c_i\leq c_m$. The scale identity is
$\prod_{j=i}^{m-1}s_j=\exp((\sigma_m-\sigma_i)/8)$.

At the top the terminal bill bound and $A_m\|x_m\|_1\leq c_m$ imply
$2C_m\|x_m\|_1\leq c_m$ and hence
$c_m\leq2(1+C_m)w$. Therefore
$\exp((\sigma_m-\sigma_i)/8)(c_i+\|x_i\|_1)
\leq2c_m\leq4(1+C_m)w$. This proves the cone clause directly from
the room and bank inequalities, on the same finite interval and history.
\end{proof}

\paragraph{Source producer retained in the lemma.}
For each step, the response is supplied by a same-history Tone record, a
paid Coupon for every keyed silent upper coefficient, and the matching
Carol/Tale witness. The Coupon fixes the input exponential entry, forced
profile, distinguished leading value, and density; its output coefficient
has the specified exact density identity and paid leading/remainder bounds.
The Carol witness includes the upper and lower forced profiles and
exponential entries at every displacement, together with common finite
support and the charge comparison used above. These are exact inputs, not
continuity limits.

The local three-card algebra which produces the paid Coupon from a
Heading/Stampede record is proved in Section~\ref{sec:three-port-history}.
It uses the common-history identities, waybill and forced-price estimate
there, together with the three selected card estimates and the finite
threshold maximum. This finite descent consumes the resulting Coupon and
Carol records; it does not reprove the card algebra. The remaining issue is
to attach those records, with their uniform prices and matching
period/degree indices, to every adjacent door of the original selected
history. That attachment is an upstream construction input here.

\subsection{Application to the terminal Cream}

We now state the calibration used by the selected field argument. Fix one
family $C$, one selected history, period $p$, terminal volume $N$, and
degree $d$. Let $Z$ be the terminal coefficient of the Cream from
Proposition~\ref{prop:ex-marked-invitation}, let $b$ be its continuous
field, and let $E>0$ be its budget. The exact terminal data are its format,
remainder bound, charge, attainment and forced-profile identity, and
density equation
\[
 Z.\mathrm{density}(U)=a_N^d R_N.\mathrm{eval}(U)b(U).
\]
Let $M$ be the Music produced on this same selected history. The paid
height record gives one scale weight sequence and a positive price
$\zeta_i$; its adjacent bounds are
\[
 0<\phi_i\leq1,\qquad
 \phi_i L_i^{d+1}\zeta_i\leq\zeta_{i+1},
 \qquad L_i=\operatorname{side}(H_i).
\]
The record also gives the local Tone/Carol witnesses, the Brine bounds,
the required endpoint compatibility, and the terminal Banks and support
data. For the lower root-test range take $k\geq \operatorname{ref}(d)$.
Write $\mathcal O(m,k)$ for the selected opening of the census at $m$ to
the working scale $k$. For each root test with $m\leq k$ and
$\mathcal O(m,k)$, the result also
requires its finite-census certificate: a finite slot set with centers at
every level from $k$ to $N$, centers transported by the scale embeddings,
the occupancy bound at $k$, and the stated Drake separation bound there.
The same certificate includes a nondecreasing occupancy schedule, the
Fiesta and Brine inequalities, the Moonlit transition kernels, and
continuity of the terminal field $b$. These are exactly the hypotheses
used by the finite-census estimate below; they are not consequences of
terminal Cream data alone.

\begin{proposition}[Same-history field calibration]
\label{prop:actual-marked-figment}
Assume $C.\mathrm{ref}(d)<N$. Under the terminal Cream and same-history
paid records above, there is one
label $v$ such that
\[
 \operatorname{puffs}(H_{N-1})
       \|Z.x-X_N^{-1}v\|_1
       \leq 2(1+C_N)Q,
 \qquad C_N=\operatorname{columns}(N),
\]
and, for every $k$ with $\operatorname{ref}(d)\leq k<N$, every $m\leq k$
with $\mathcal O(m,k)$, every motion
$s$, and every root $o$ at scale $N$ supplied with the stated finite-census
certificate,
\[
 \operatorname{rootTest}_{N,o}\bigl(s\cdot(b-C.O(v))\bigr)
       \leq \operatorname{cider}(d,m,k)\,E\,
               \operatorname{flight}_{-\gamma}(N,p).
\]
Here $X_N$ is the terminal coefficient calibration isomorphism on the
selected family, and $C.O(v)$ is its zero-source local observable.
The label $v$ and all bounds use the same $C$, history, and Music.
\end{proposition}
\begin{proof}
The selected paid record first supplies the Music $M$ and its Heard and
Responds fields from the bottom reference $r=\operatorname{ref}(d)$. Its adjacent Tone and Carol
records, together with the exact terminal Cream bounds, are the inputs of
Lemma~\ref{lem:actual-marked-antilles} with $m=N$ and terminal cone
parameter $w=Q$, as specified by the Cream charge bound. The larger budget
$E$ is reserved for the root-test error estimate. This produces one
backward sequence on the original history. Let $a$ be the primary of its
state at $r$, and set $v=X_N(a)$. By invertibility,
$X_N^{-1}v=a$.

The terminal coordinate identity gives
$(Q_N).bar_M(-a).x=Z.x-X_N^{-1}v$. The cone estimate at level $N$ and the
terminal spread bound give $Q_N.cost_M(-a)\leq2(1+C_N)Q$. The last step
estimate from scale $N-1$ says
\[
 \operatorname{puffs}(H_{N-1})\|(Q_N).bar_M(-a).x\|_1
       \leq Q_N.cost_M(-a).
\]
Combining these inequalities proves the first displayed estimate. Its
hypothesis $\operatorname{ref}(d)\leq N-1$ follows from the explicit
assumption $\operatorname{ref}(d)<N$.

For the root-test estimate choose
\[
 w=\frac{\zeta_k/4}
          {\|\operatorname{alto}(H,k,d)\|\,\operatorname{escrow}(N,k,E)}.
\]
Every factor in the denominator is positive. The height record gives
$\zeta_k>0$, the physical scale gives
$\|\operatorname{alto}(H,k,d)\|=a_k^{-d}>0$, and $E>0$ gives positive
escrow. Marched weights are $w_i=w\prod_{q=k}^{i-1}\phi_q$. The exact
recurrence for the physical scale and the gain is
\[
 \|\operatorname{alto}(H,i+1,d)\|
       =L_i^d\|\operatorname{alto}(H,i,d)\|,
 \qquad
 \operatorname{gain}(N,i+1)
       =\operatorname{gain}(N,i)/\operatorname{puffs}(H_i).
\]
Since $1\leq\operatorname{puffs}(H_i)\leq L_i$, the adjacent paid bound
implies inductively
\[
 w_i\|\operatorname{alto}(H,i,d)\|
       \frac{E}{\operatorname{gain}(N,i)}\leq \frac{\zeta_i}{4}
       \qquad (i\geq k).
\]
The base case is equality at $k$ by the definition of $w$. For the step,
the left side is multiplied by at most $\phi_iL_i^{d+1}$; the displayed
price recurrence pays this by $\zeta_{i+1}/\zeta_i$. Thus the exact
march/height inequality required for the finite root estimate holds at
every later scale.

For precision, these are the fields of the finite-copy construction (the
source's Yeast record). At level $i$, the residual $Q_i.\mathrm{sil}$ is
keyed; the coefficient used by the Yeast record is the translated voice
$Z_i=(Q_i).\mathrm{bar}_M(-a)$ with its marched bill. The `keys` identity
transfers format and the remainder bound to this translated voice, which
need not have zero leading coordinate. The height bound on
$\mathrm{gain}(N,i)(\mathrm{cost}_i+\|Z_i.x\|)$, combined with the
per-level price inequality for
\[
 \frac{w_i\|\operatorname{alto}(H,i,d)\|E}{\mathrm{gain}(N,i)},
\]
gives the Yeast rate clause by multiplication. The
Coupon's density debt is exactly the density transition in the Yeast vote;
its Carol witness supplies the accompanying finite Tale support and the
forced-profile/exponential-entry comparison. At the terminal level let
$c=(M.\mathrm{tone}\,N).\mathrm{index}(-a)$. The top clause of `Sung` gives
$(M.\mathrm{tone}\,N).\mathrm{index}=\mathrm{id}$, so $c=-a$; the
terminal spread bound gives $\|c\|_1\leq2E$. Put
$j=\|\operatorname{alto}(H,N,d)\|$ and let $w_N$ be the terminal price.
The height estimate $w_NjE\leq\zeta_N/4$ implies
\[
 2w_Nj\|c\|_1\leq\zeta_N.
\]
The terminal `erected` estimate, applied at this response cost, supplies an
added price $A_{\rm ins}$ with $w_NjA_{\rm ins}\leq u_N/2$. If
$B_{\rm cream}$ is the Cream's original attainment price, the Bank
hypothesis
\[ B_{\rm cream}\zeta_N\leq u_NE, \]
together with
$w_NjE\leq\zeta_N/4$, gives
$w_NjB_{\rm cream}\leq u_N/4\leq u_N/2$. Hence the combined price
$A_{\rm ins}+B_{\rm cream}$ fits the Yeast allowance $u_N$. The forced profile is
preserved by the Cream hook identity and linearity of
the played forced-profile map. Finally, the Yeast bill is the Cream density
equation after the terminal `originate' identity: the two recentering
terms have density equal to the Music readout, leaving precisely the
terminal marked density. These are all same-history records; no new
terminal volume or history is selected in assembling the Yeast.

The Yeast residual coefficients are keyed. The chosen negative bottom
primary makes the bottom translated voice have zero leading coordinate.
The Accord identities give the Wilson push and all-displacement profile
updates; the Cream gives the
terminal forced hook and density identity; and the Bank and footprint
cover pay the forced correction. The output is the exact Yeast identity
for the difference between $b$ and the original-history observable with
primary $-a$. Since the terminal calibration is linear,
\[
 \operatorname{initial}(-a)= -C.O(X_Na)=-C.O(v),
\]
so that difference is $b-C.O(v)$. To spell out the finite root estimate,
fix an opened census at scale $m$ and write $K=|\mathrm{Site}(m,p)|$ for
its number of slots. The selected root-test data must also give the exact
inputs used in the finite census lemma: the migration of the census
centers from scale $k$ through $N$, the occupancy and Drake separation
bounds at scale $k$, the monotone occupancy schedule, the Fiesta and
Brine inequalities, the Moonlit transition kernels, and continuity of
$b$. With constant positive slot weight $w$, the finite census estimate
(obtained by applying the paid transition bound to the product over all
$K$ slots) is
\[
 \|\Omega\|w^K\leq
   \exp\!\left(|\mathrm{Site}(k,p)|\,\mathrm{Treasury}_k+5K\right),
 \qquad
 \Omega=\int\prod_{j\in\mathrm{Site}(m,p)}U_j\,d\mu_{N,p}.
\]
Here $U_j$ is the $j$th translated/reflected census copy, with the chosen
motion applied, of the residual field $b-C.O(v)$. Thus $\Omega$ is exactly
the moment in the definition of the displayed root test.
For each slot, put $x_k=Z_k.x$ for the leading coefficient of its Yeast
voice; its remainder allowance is $\phi_k$, and the marched weight has
$w_k=w$. The rate bound at the opened scale $k$ is
\[
 w_k\|\operatorname{alto}(H,k,d)\|C_k
       (\phi_k+\|x_k\|_1)\leq1,
 \qquad C_k=\operatorname{columns}(k)\geq1.
\]
The local hull estimate bounds the magnitude of its exponential entry by
$\|x_k\|_1+2C_k\phi_k$. Since
$\|x_k\|_1+2C_k\phi_k\leq2C_k(\phi_k+\|x_k\|_1)$, the weighted entry price is
at most $2$. The forced-profile estimate gives weighted total forced price
at most $1$. The action allowance is $f=1$. Thus the marked estimate of
Section~\ref{ex-centered-sail} charges $x+y+2f=5$ per slot (the more
general choice $x=3,y=1$ gives the $6$-unit bound). The constant $5K$ is
exactly the sum of the per-slot tolls, and the extensive term is the cost
of all scale-$k$ sites. The aligned side-count identity is
\[
 |\mathrm{Site}(k,p)|=K(a_m/a_k)^4.
\]
Consequently, with
$g=(a_m/a_k)^4\mathrm{Treasury}_k+5$, we have
$\|\Omega\|w^K\leq e^{Kg}$. If $\Omega=0$ the desired inequality is
immediate; otherwise the same inequality still gives
$\|\Omega\|^{1/K}\leq e^g/w$. By definition the normalized census root
test is $\|\Omega\|^{1/K}$, hence
\[
 \operatorname{rootTest}_{N,o}(s\cdot(b-C.O(v)))
       \leq \frac{\operatorname{stakes}(m,k)}{w}.
\]
Finally, the definitions of escrow, flight, and the stake convert this
ratio exactly:
\[
 \frac{\operatorname{stakes}(m,k)}{w}
 =\operatorname{stakes}(m,k)
   \frac{4\|\operatorname{alto}(H,k,d)\|}{\zeta_k}
   \operatorname{escrow}(N,k,E)
 =\operatorname{cider}(d,m,k)E\operatorname{flight}_{-\gamma}(N,p).
\]
This is uniform over the stated tests when the height, transition, and
Cream records, including the listed census hypotheses, are uniform over
them. No subsequence or replacement history enters the construction.
\end{proof}

\paragraph{Boundary of this producer.}
The proof above expands the terminal recentering, finite backward induction,
cost conversion, marched-weight calculation, and the final use of the
finite-copy estimate. It does not derive the selected history's Tone/Coupon
producer, its all-volume height ledger, or the selected-history binding of
the terminal Cream. Those exact suppliers are listed separately in the
evidence record; none follows merely from the existence of a portfolio.
The precise remaining certificate, on one fixed family and one fixed field
of adjacent gauge doors, is as follows. For each degree $d$, terminal level
$N\geq\operatorname{ref}(d)$, and period $p$, it must provide a Music whose
terminal voice is the chosen factory representative, whose played fields
realize the factory entries, which is heard at every level from
$\operatorname{ref}(d)$ to $N$, and which responds across each intervening
door. The same Music must have Carol witnesses at all those doors and the
terminal forced-profile hook. For every $i\geq\operatorname{ref}(d)$ it
must also provide the paid local Coupon and its all-displacement Carol/Tale
comparison, with the uniform height, room, bank, Brine and monotone
occupancy prices used above. These are same-family, same-history
quantifiers: choosing a different Music, history, or door field for each
period or terminal volume is insufficient.

In the source construction this record is attached by the selected
funded-history route: its history package gives the nested selected field
and adjacent climbing doors; its finite Music package supplies the terminal
factory realization and full-height response history; the local three-card
Heading/Stampede calculation of Section~\ref{sec:three-port-history}
prices each Coupon; and the all-level service ledger supplies the Geyser
height and transition bounds. The remaining proof obligation is to derive
these inputs uniformly from the same original funded history, including the
terminal Cream attachment. The local three-card calculation, the
conditional calibration proved here, and the finite-census inequality do
not by themselves establish that attachment.
